% !PRINTED 37
\weberpage{37}
%<<CONT

\begin{equation}
A=QB+C,
\tag{13}
\end{equation}
\begin{equation}
 C=c_0x+c_1,
\tag{14}
\end{equation}
worin
\begin{equation}
 c_0={6a_0a_2-2a_1^2\over 9a_0},\qquad
 c_1={9a_0a_3-a_1a_2\over 9a_0}.
\tag{15}
\end{equation}
Wenn $c_0$ gleich Null ist, so ist hiermit der Algorithmus schon geschlossen; wenn $c_1$ von Null verschieden ist, dann haben $A$ und $B$ keinen gemeinsamen Theiler; ist aber $c_1$ auch gleich Null, so ist $B$ selbst der grösste gemeinsame Theiler von $A$ und $B$, d. h. $A$ ist durch $B$ theilbar. Die Bedingungen hierfür sind also:
\begin{equation}
 3a_0a_2-a_1^2=0,
 \qquad
 9a_0a_3-a_1a_2=0.
\tag{16}
\end{equation}
Ist $c_0$ nicht gleich Null, so gehen wir einen Schritt weiter und setzen:
\begin{equation}
 B=PC+D,
\tag{17}
\end{equation}
worin $D$ constant wird und den Ausdruck erhält:
\begin{equation}
 D={a_2c_0^2-2a_1c_0c_1+3a_0c_1^2\over c_0^2}.
\tag{18}
\end{equation}
Ist dieser Ausdruck von Null verschieden, so sind $A$ und $B$ \gesperrt{theilerfremd}, ist er gleich Null, so haben $A$ und $B$ den grössten gemeinschaftlichen Theiler $C$. Setzen wir für $c_0,c_1$ die Werthe (15) ein, lassen den Nenner weg, heben noch den von Null verschiedenen Factor $9a_0$ heraus und kehren das Vorzeichen um, so erhält diese Bedingung nach einfacher Rechnung die Gestalt:
\begin{equation}
 a_1^2a_2^2+18a_0a_1a_2a_3-4a_0a_2^3-4a_1^3a_3-27a_0^2a_3^2=0.
\tag{19}
\end{equation}
Sie ist, wie man leicht durch Rechnung oder auch aus (18) sieht, auch dann erfüllt, wenn die Bedingungen (16) bestehen, und ist also die nothwendige und hinreichende Bedingung dafür, dass $f(x)$ und $f'(x)$ einen gemeinsamen Theiler haben. Die linke Seite von (19), die eine ganze rationale und homogene Function der Coëfficienten von $f(x)$ ist, heisst die \gesperrt{Discriminante der Function $f(x)$}.

Wir leiten noch einen Satz aus dem Algorithmus (3) her, der oft angewandt wird.

Aus der ersten dieser Gleichungen folgt:
\begin{equation}
 A''=A-Q'A',
\tag{20}
\end{equation}
und wenn man diesen Werth von $A''$ in die folgende Gleichung einsetzt:
\[
 A'''=(1+Q'Q'')A'-Q''A,
\]
%CONT>>
